% ECONOMICS_PROBLEM: {"id": "C23-251", "title": "Causal effects of treatment histories in panel data", "jel_code": "C23", "source_id": "OP-0599"}
% !TeX program = xelatex
\documentclass[11pt,a4paper]{article}
\usepackage[margin=23mm]{geometry}
\usepackage{fontspec}
\setmainfont{Latin Modern Roman}
\usepackage{amsmath,amssymb,mathtools}
\usepackage{xcolor,enumitem,booktabs,longtable,array,xurl,hyperref}
\definecolor{muted}{HTML}{53636C}
\hypersetup{colorlinks=true,urlcolor=blue,linkcolor=blue}
\setlength{\parindent}{0pt}
\setlength{\parskip}{5pt}
\newcommand{\R}{\mathbb R}
\newcommand{\E}{\mathbb E}
\newcommand{\N}{\mathbb N}
\newcommand{\ind}{\mathbf 1}
\newcommand{\Var}{\operatorname{Var}}
\newcommand{\Cov}{\operatorname{Cov}}
\newcommand{\supp}{\operatorname{supp}}
\DeclareMathOperator*{\argmin}{arg\,min}
\DeclareMathOperator*{\argmax}{arg\,max}
\newcommand{\entrymeta}[3]{{\sffamily\small\color{muted}\textbf{\texttt{#1}}\quad|\quad #2\par #3\par}}
\newcommand{\Problemref}[1]{\texttt{#1}}
\newcommand{\JELref}[2]{\texttt{#2} (JEL \texttt{#1})}
\begin{document}
\section*{Causal effects of treatment histories in panel data}
\textbf{Problem ID:} \texttt{C23-251}\quad \textbf{JEL:} \texttt{C23}\par
% BEGIN ECONOMICS STATEMENT
\label{problem:OP-0599}
% ECONOMICS_PROBLEM: {"local_id": "ECON-014", "title": "Causal effects of treatment histories in panel data", "kind": "R", "audit_date": "2026-10-08", "source_identity_checked": true, "agenda_or_question_support_checked": true, "source_role": "see per-entry audit and locators", "current_open_certified": false}


\label{audited:ECON-014}

{\sffamily\small\color{muted}\textbf{ECON-014} | R | Source-backed agenda; exact formalization is editorial\par}

\subsection*{Definitions and formal target}

Observe $(Y_{it},D_{it},X_{it})$ for units $i$ and dates $t\leq T$, with binary treatment $D_{it}$. Let $Y_{it}(d_{1:t})$ denote a potential outcome under the complete treatment history; no anticipation means future treatments do not affect it. Latent unit factors and time-varying shocks may affect both treatment and outcomes.

For specified histories $d_{1:t},d'_{1:t}$ and population $G$, the target is
\[
 \tau_t(d,d';G)=\mathbb E[Y_{it}(d_{1:t})-Y_{it}(d'_{1:t})\mid i\in G].
\]
Characterize restrictions on selection, carryover, and latent factors under which this target is identified from the observed panel law, and obtain valid inference when the number of observed histories grows with $T$.

\subsection*{Clarifications and required assumptions}

All history-indexed outcomes are defined on one probability space. Consistency is $Y_{it}=Y_{it}(D_{i1:t})$; the main individual-history notation additionally assumes no interference between units. The group $G$ is fixed by baseline covariates, with positive population probability, rather than selected using realized treatment or outcomes. An observed-law class alone does not define unobserved potential outcomes: use structural models $m\in\mathcal M$ and target $\tau_t(m;d,d',G)$. The question asks for restrictions making this target constant across models with the same observed panel law.

\subsection*{Variants and changes of scope}

\textbf{Finite carryover.} Fix $q$ and require $Y_{it}(d_{1:t})$ to depend only on $d_{\max(1,t-q):t}$; compare with unrestricted histories. \textbf{Sequential exchangeability.} At each date, impose independence of assignment and specified future potential outcomes conditional on observed past, with positive conditional treatment probabilities. \textbf{Latent factors.} Fix a rank bound and restrictions connecting untreated outcome factors to treatment assignment; arbitrary hidden confounding is not resolved by writing a factor equation. \textbf{Asymptotics.} State whether $n\to\infty$, $T\to\infty$, or both, and how the tested history set grows.

\subsection*{Audit conclusion and source relationship}

Section 10.1 of the checked author version supplies the dynamic-outcomes agenda. The history contrast and latent-factor extensions are editorial specifications. Existing sequential-treatment methods are not claimed nonexistent.

\subsection*{References and exact source roles}

{\footnotesize\raggedright
\par\noindent\textbf{[S1]} Dmitry Arkhangelsky and Guido Imbens (2024). \emph{Causal Models for Longitudinal and Panel Data: A Survey}. The Econometrics Journal 27(3), C1--C61, DOI: 10.1093/ectj/utae014. Checked author version: arXiv:2311.15458v3, 25 June 2024.\par \textit{Verified locator / source role:} Section 10.1 in the checked June 2024 author version. The journal landing page was verified for metadata; the readable author PDF was used for the agenda.\par \url{https://arxiv.org/pdf/2311.15458}\par\smallskip
}

\subsection*{Common conventions: Additional-source status and mathematical conventions}

\label{audited:conventions}

Of the 100 supplied ECON entries, 65 distinct problems appear here; 32 close
matches refine earlier entries, and three duplicate questions map to existing
entries. Appendix~\ref{app:audited-crosswalk} lists every destination.
The attachment's P/R/S/U distinctions and source audits are retained. Its
precise mathematical formalizations are editorial unless the individual audit
says otherwise. Candidate subsidy targets ECON-004--006 retain their U flags;
the resolved approval-core question ECON-007 updates OP-0202 above.
The following supplied conventions govern both this part and the attributed
ECON refinements elsewhere in the catalogue, subject to local restrictions.

\textbf{Parameterized questions.} A problem family lists inputs that must be fixed before an instance is evaluated: population, horizon, policies, information, data law, model class and welfare weights. These are required choices, not quantities the citations are assumed to specify. Undefined applications should not be treated as identified estimands.

\textbf{Indices and integrability.} Sets of agents, firms and regions are finite unless a population measure is explicitly used. \(i,j\) index units, \(r\) regions, \(t\) dates and \(h\) horizons. \(\mathbb E\) and \(\Pr\) refer to the declared population and law; \(\mathbf1\{A\}\) is an indicator. Every displayed expectation and discounted sum needs existence in the stated domain. Infinite horizons use a discount in \((0,1)\) and a convergence condition; finite horizons may use discount one.

\textbf{Optimization and existence.} A displayed \(\max\), \(\min\), or \(\arg\max\) is conditional on attainment. For finite feasible menus this is immediate; general spaces need continuity and compactness/coercivity or another existence argument. Without attainment, the target is the supremum/infimum and arbitrarily near-optimal feasible choices. \(\inf\varnothing=+\infty\). Empty feasibility and an unidentified target are different issues.

\textbf{Potential outcomes.} \(Y_i(z)\) is the outcome under a specified intervention, not the observed conditional mean among people choosing \(z\). In binary comparisons \(\Delta Y_i=Y_i(1)-Y_i(0)\). Specify assignment, treatment versions, compliance, information and observation horizon. Interference is allowed under a full policy vector; compressed exposure notation needs a justified mapping. No interference, consistency and overlap are substantive assumptions, not automatic consequences of notation.

\textbf{Populations and observation.} Fix baseline eligibility and population weights. Conditioning on post-policy employment, survival, relocation or program uptake can change the population being compared. Death is not ordinary loss to follow-up; outcomes truncated by death need a composite convention, joint survival target, or explicitly defined survivor stratum. Fertility-changing interventions need a population functional rather than an unexplained fixed descendant cohort. Measurement and missingness models must be supplied.

\textbf{Identification.} A structural model \(m\in\mathcal M\) implies observable law \(P_m\) and target \(\theta(m)\). Identification means
\[
 P_{m_1}=P_{m_2}\ \Longrightarrow\ \theta(m_1)=\theta(m_2)
 \quad\text{for all }m_1,m_2\in\mathcal M .
\]
Without identification, the target is
\[
 \Theta(P;\mathcal M)=\{\theta(m):m\in\mathcal M,\ P_m=P\}.
\]
Writing \(\theta(P)\) before establishing identification can hide incompatible latent models. Confidence sets additionally need a sampling law, confidence level, asymptotic sequence and uniformity class. Point identification, coverage and informative precision are separate properties.

\textbf{Mediation.} Total effects of randomized assignment are distinct from cross-world natural indirect effects. Belief/effort-channel effects require a meaningful mediator intervention and additional measurement, exclusion or mediation assumptions. Randomization of the initial policy alone does not supply those assumptions. Report bounds or interventional alternatives when the stronger target is unsupported.

\textbf{Equilibrium.} A counterfactual \(W(z)\), \(p(z)\), or \(\mu_t^z\) requires individual optimization, feasibility, government/resource budgets, market clearing and state-transition laws. State equilibrium existence and selection, or report ranges over compatible equilibria. Initial-state integration includes subsequent endogenous transitions; current-state integration must use the policy-dependent distribution. Reduced-form invariance after policy is not assumed.

\textbf{Welfare accounts.} Scores, utility, health, dollars and ecological quantities require explicit conversion before addition. Fees, taxes, subsidies, wages and extortion payments are transfers between parties; distributional weights can make them welfare relevant, but their face value is not automatically a resource loss. State ownership, geographic boundaries, foreign-income weights and fiscal finance. Do not add profits to household welfare already including dividends, or count land capitalization and the underlying benefit twice. A benefit-cost expression must distinguish gross benefits from net welfare and subtract every real cost exactly once.

\textbf{Derivatives and risk.} Log derivatives require positive arguments; local responses require admissible perturbations and specified equilibrium branches. Differentiation under expectations needs regularity/domination, and boundaries may allow only one-sided derivatives. Risk uses a specified law; ambiguity uses a declared nonempty law set on a common history space. Adaptive policies must be nonanticipative. A robust criterion is an editorial choice unless attributed explicitly.

\textbf{Scope overlaps.} Allocation, collective choice, contracts, household bargaining and coordinated policy overlap economics with optimization and game theory. They are retained because they appeared in the original catalogue; no claim of a strictly game-theory-free selection is made.
% END ECONOMICS STATEMENT
\end{document}
